% Accuracy-oriented extensions (scope of 2026-11-18: accuracy first, inference cost does not matter).
% Skeleton of the EPFL report's new results; numbers from TASKS.md §9 (2026-11-20). Every cell
% marked \pending is a run that has not finished; \todo marks text to write after the decision
% gate of 2026-12-04. All AR values: BOP19 average recall, LM-O / YCB-V / T-LESS test with the
% BOP default detections. Validation: LM-O real validation images (the 200 BOP19 images of each
% object), YCB-V four held-out real videos, T-LESS held-out real rings, all with noisy boxes
% (each side moved by up to 10\,\% of the box size).
\section{Accuracy-Oriented Extensions}
\label{sec:accuracy}

The main model of the previous sections trades accuracy for speed. This section drops that
constraint: inference cost and model size do not matter, only the BOP19 AR does. Every change is
chosen on the validation images and adopted only if the mean of two seeds beats the baseline's
mean on both the noisy-box validation score and the test AR, by any margin
(\cref{sec:acc:rule}). Results trained on PBR images only and results that also use real training
images are always reported separately.
\todo{Rewrite the abstract, introduction and conclusion for this scope after the 2026-12-04 decision gate.}

\subsection{Selection Rule}
\label{sec:acc:rule}
Each candidate is trained with two seeds and compared with the baseline's two seeds under the same
budget. The validation score is $(\mathrm{AR}_\mathrm{MSSD}+\mathrm{AR}_\mathrm{MSPD})/2$ on real
validation images with perturbed boxes, so that the choice does not use the test set and sees the
box errors of a detector; the test AR is read once per candidate. A change is adopted when both
means improve.

\subsection{Backbone and Input Resolution}
\label{sec:acc:backbone}
\begin{table}[t]
  \centering
  \small
  \begin{tabular}{@{}llcccc@{}}
    \toprule
    Backbone & Input & Val. a / b & Val. mean & Test a / b & Test mean \\
    \midrule
    ViT-B/16 (baseline) & 384 & .885 / .880 & .883 & .664 / .658 & .661 \\
    ViT-B/16 & 512 & .881 / .884 & .882 & .657 / .664 & .661 \\
    ViT-B/16 & 640 & \pending{running} & & & \\
    ViT-L/16 & 512 & \pending{seed a} / .890 & & \pending{seed a} / .681 & \\
    ViT-H/14 & 518 & \pending{running} & & & \\
    \bottomrule
  \end{tabular}
  \caption{\textbf{Backbone and input size on LM-O} (PBR only, 2$\times$ budget, two seeds).
  Validation: noisy-box score on the real validation images; test: BOP19 AR with detections.}
  \label{tab:acc_backbone}
\end{table}
A larger input alone does not help (ViT-B at 512\,px equals the 384\,px baseline on both numbers,
\cref{tab:acc_backbone}); LM-O crops are usually smaller than 384\,px natively, so more pixels add
no detail. The first ViT-L/16 seed is clearly better (+0.008 validation, +0.021 test) and also
much better on the synthetic validation images (0.942 vs.\ 0.909).
\todo{Decide with the second seed; ViT-B/640 and ViT-H/518 are still running.}

\subsection{Training Budget}
\label{sec:acc:budget}
Going from $2\times$ to $4\times$ the number of samples per epoch gave +0.011 AR. At $8\times$
the first seed reaches 0.673 test AR and 0.886 noisy-box validation, against 0.668 / 0.886 at
$4\times$ (same recipe, first seeds); the second seed is \pending{running} (\cref{tab:acc_budget}).
\begin{table}[t]
  \centering
  \small
  \begin{tabular}{@{}lccc@{}}
    \toprule
    Budget & Val. (a / b) & Test (a / b) & Mean test \\
    \midrule
    $4\times$ & .886 / .880 & .668 / .664 & .666 \\
    $8\times$ & .886 / \pending{running} & .673 / \pending{running} & \\
    \bottomrule
  \end{tabular}
  \caption{\textbf{Training budget on LM-O} (PBR only, baseline recipe).}
  \label{tab:acc_budget}
\end{table}

\subsection{One Network per Object}
\label{sec:acc:per_object}
\begin{table}[t]
  \centering
  \small
  \begin{tabular}{@{}lcccc@{}}
    \toprule
    & \multicolumn{2}{c}{One network (shared)} & \multicolumn{2}{c}{One per object} \\
    Dataset & Val. & Test & Val. & Test \\
    \midrule
    LM-O  & .883 & .666 & .879 & .669 \\
    YCB-V & .691 & .698 & .674 & .713 \\
    T-LESS & .878 & .762 & .845 & .736 \\
    \bottomrule
  \end{tabular}
  \caption{\textbf{Shared network vs.\ one network per object} (PBR only, $4\times$ budget per
  object, mean of two seeds). Per-object models are not adopted on any dataset: they are better on
  the test set of LM-O and YCB-V but worse on validation, and clearly worse on T-LESS.}
  \label{tab:acc_per_object}
\end{table}
Eight to thirty separate models cost up to thirty times the weights and do not pay off by the
selection rule (\cref{tab:acc_per_object}). On YCB-V the test AR improves by 0.015 but the
validation score drops by 0.017, so the improvement is not trusted.

\subsection{Training with Real Images}
\label{sec:acc:real}
\begin{table}[t]
  \centering
  \small
  \begin{tabular}{@{}lccccc@{}}
    \toprule
    & \multicolumn{2}{c}{PBR only} & \multicolumn{2}{c}{PBR + real} & \\
    Dataset & Val. & Test & Val. & Test & Real data \\
    \midrule
    LM-O  & .883 & .661 & .972\textsuperscript{*} & .681 & LM, no BOP19 images \\
    YCB-V & .676 & .685 & .752 & .764 & \texttt{train\_real} \\
    T-LESS & \pending{running} & & & & \texttt{train\_primesense} \\
    \bottomrule
  \end{tabular}
  \caption{\textbf{Adding the real training images} ($2\times$ budget, mean of two seeds; the two
  groups of columns are different leaderboard categories). \textsuperscript{*}The LM-O real
  validation images lie about $5^\circ$ from real training views, which inflates the validation
  score of the real-trained model; its test AR is not affected.}
  \label{tab:acc_real}
\end{table}
Real images help on both datasets (\cref{tab:acc_real}): $+0.079$ test AR on YCB-V and $+0.020$ on
LM-O. On YCB-V the validation videos are held out entirely, so the validation gain
($+0.076$) is a fair estimate. These models belong to the PBR + real category of the BOP
leaderboard and are reported as such.

\subsection{Dense Heads and PnP}
\label{sec:acc:dense}
\begin{table}[t]
  \centering
  \small
  \begin{tabular}{@{}lcccc@{}}
    \toprule
    Pose mode & Val. & Test & Time (ms) & Fallbacks \\
    \midrule
    Direct regression, no dense heads (baseline) & & & & \\
    Direct regression, with dense heads & \pending{} & \pending{} & & \\
    PnP on mask + XYZ & \pending{} & \pending{} & & \\
    PnP from the regressed pose & \pending{} & \pending{} & & \\
    Regressed rotation + PnP translation & \pending{} & \pending{} & & \\
    \bottomrule
  \end{tabular}
  \caption{\textbf{Multi-head network on LM-O} (PBR only, $2\times$ budget). One ViT predicts the
  pose directly and, from its patch tokens, a visible-object mask and normalised object
  coordinates; PnP-RANSAC, a refinement from the regressed pose, or a translation solved for the
  regressed rotation give the final pose. ``Fallbacks'': instances where PnP failed and the
  regressed pose was kept.}
  \label{tab:acc_dense}
\end{table}
\todo{Describe the decoder (\texttt{DenseHead}), the targets (depth + ground-truth pose + visible
mask, normalised by the model box), the losses and the symmetric-object handling once the six
\texttt{lmo-m18-*} runs are compared; weights 1 / 5 / 0.2 for the XYZ term are the candidates.}

\subsection{Render-and-Compare Refinement}
\label{sec:acc:refine}
\begin{table}[t]
  \centering
  \small
  \begin{tabular}{@{}lccccc@{}}
    \toprule
    Refinement iterations & 0 & 1 & 2 & 3 & 4 \\
    \midrule
    LM-O  & .669 & .616 & .584 & .567 & .557 \\
    YCB-V & .694 & .693 & .694 & .689 & .689 \\
    T-LESS & .762 & .743 & .731 & .722 & .719 \\
    LM-O, our refiner & .669 & \pending{} & \pending{} & \pending{} & \pending{} \\
    \bottomrule
  \end{tabular}
  \caption{\textbf{CosyPose's PBR-only refiner~\cite{labbe2020cosypose} on our poses} (test AR
  with detections, PBR-trained refiner and no further training; 0 = our poses as they are).}
  \label{tab:acc_refine}
\end{table}
The public refiner brings no gain (\cref{tab:acc_refine}): it corrects large errors (on YCB-V, a
noisy ground truth of 15$^\circ$ / 40\,mm goes down to 4.8$^\circ$ / 21\,mm in two iterations) but not
errors below about $3^\circ$ and 15\,mm, which is where our poses already are, and it makes LM-O
and T-LESS worse. Each iteration costs about 0.04\,s per image.
\todo{Decide whether to train our own refiner on perturbations of our poses (M19 R2).}
